\section{匹配系数}
在固定 $t>0$ 流时间处完成强子矩阵元的连续极限后，
与 $t=0$ 物理重正化矩阵元的联系由 SFTX 给出~\cite{Luscher:2013vga}。
twist-2 算符 $O_n(x,t)$ 的 SFTX 一般包含形如
$1/t^{m} \delta_{\mu_i \mu_j}$ 的幂次发散项，
其中 $m=(d_n - d_p)/2$ 依赖于 $O_n$ 的量纲 $d_n$
以及低维算符的量纲 $d_p$。不过这些项按连续 O($4$) 对称性分类，
因此对称化无迹算符 $\widehat{O}_n$ 的 SFTX
只接收来自 $t=0$ 相应无迹算符的贡献。

我们采用上一节引入的环场重正化带流费米子场，并定义环化 twist-2 算符
\bea
\widehat{\rO}_n^{rs}(x,t) &=&
\rchibar^r(x,t) \gmuopen1 \lrDmu2 \cdots \lrDmuclose{n} \rchi^s(x,t) - \nonumber \\
&-& \text{含~}\delta_{\mu_i \mu_j}\text{~的项}\,,
\label{eq:flowed_t2_ringed}
\eea
其中的减除保证所得算符迹为零。$\widehat{\rO}_n^{rs}$ 的 SFTX 为
\be
\widehat{\rO}_n^{rs}(t) = c_n(t,\mu)\widehat{O}_{n,\MS}^{rs}(\mu) + \cdots\,,
\ee
这里只显式写出对流时间与重正化标度的依赖；被略去的项是 $t$ 正幂次的高阶贡献。

匹配系数通过考虑含有带流算符 $\widehat{\rO}_n^{rs}(t)$ 的
离壳裁剪单粒子不可约（1PI）格林函数确定。
无质量费米子的匹配方程
\be
\left\langle \psi^r \widehat{\rO}_n^{rs}(t) \psibar^s \right\rangle = c_n(t,\mu)
\left\langle \psi^r \widehat{O}_{n,\MS}^{rs}(t=0,\mu) \psibar^s \right\rangle
\label{eq:matching_eq}
\ee
在 $D = 4 - 2 \epsilon$ 中于微扰论单圈求解，取
\be
c_n(t,\mu) =  1 + \frac{\gbar^2(\mu)}{\left(4 \pi\right)^2}c_n^{(1)}(t,\mu) +
O(\gbar^4)\,.
\label{eq:cn_1l}
\ee
匹配系数 $c_n(t,\mu)$ 依赖重正化耦合，
并与计算匹配系数时的软标度无关。
我们在补充材料~\cite{SupplMat}中描述其计算，
并基于它们满足重正化群方程的观察，给出次领头对数（NLL）的重求和表达式。

匹配系数的最终结果为
\be
c_n^{(1)}(t,\mu) = C_F \left[ \gamma_n \log \left(8 \pi \mu^2 t \right) + B_n\right]\,,
\label{eq:matching_1l}
\ee
其中 $\gamma_n = 1 + 4 \sum_{j=2}^n \frac{1}{j} - \frac{2}{n(n+1)}$，而
\bea
B_n &=& \frac{4}{n(n+1)} + 4 \frac{n-1}{n}\log 2 + \frac{2-4 n^2}{n(n+1)}\gamma_E - \nonumber \\
&-& \frac{2}{n(n+1)}\psi(n+2) + \frac{4}{n}\psi(n+1) - 4 \psi(2) - \nonumber \\
&-& 4 \sum_{j=2}^n \frac{1}{j(j-1)} \frac{1}{2^j} \phi(1/2,1,j) - \log \left(432\right)\,.
\label{eq:finite}
\eea
式中 $\gamma_E = 0.57221\ldots$ 是 Euler 常数，$\psi(z)$
为双伽马函数，$\phi(z,s,a)$ 是 Lerch 超越函数，定义为
$\phi(z,s,a) = \sum_{k=0}^{\infty} \frac{z^k}{(k+a)^s}$。

$\gamma_n$ 的表达式与 twist-2 算符的单圈反常量纲一致~\cite{Gross:1974cs}，
为计算提供了受欢迎的检验。
对 $n>2$，式~\eqref{eq:finite} 中的有限部分 $B_n$ 是新结果。
$n=2$ 的结果此前已在文献~\cite{Makino:2014taa} 分析能量—动量张量时得到，
我们重现了其结果。

带流矩阵元的匹配如此简单这一事实，
开启了计算任意阶部分子分布函数矩的可能。
一个可操作流程如下：先给定 $n$，
构造对称无迹算符 $\widehat{O}_n^{rs}$；
然后用基于谱分解的标准技术计算其在强子态间的矩阵元。
要在三点的格点计算中运用谱分解，
必须保证强子态内插场与带流算符之间的物理距离
远大于流时间半径 $\sqrt{8t}$。完成连续极限后，
$\MS$ 方案中的重正化矩阵元简单地由下式得到
\be
A_n^{\MS}(\mu) =
c_n(t,\mu)^{-1} A_n(t) \,.
\label{eq:x_MS}
\ee
只有当带流场定义为环场时该匹配才正确；
若带流场在其他方案中重正化，则匹配系数的有限部分须作修改。

如果匹配成功，式~\eqref{eq:x_MS} 左端应与流时间无关。
偏离来自以 $t$ 线性项出现的更高维算符，
以及匹配系数微扰展开的高阶项。
为缓解这些可能的系统误差，重要的是把本文的计算推广到 O($\gbar^4$) 阶，
并对匹配后强子矩阵元的残余流时间依赖进行透彻的数值研究。
